\documentclass[11pt]{article}
\usepackage[utf8]{inputenc}
\usepackage{amsmath,amssymb,amsthm}
\usepackage{hyperref}
\usepackage{listings}
\usepackage{xcolor}
\usepackage[margin=1in]{geometry}

\lstset{
    basicstyle=\ttfamily\small,
    breaklines=true,
    frame=single,
    backgroundcolor=\color{gray!5},
    numbers=none,
    xleftmargin=4pt,
    xrightmargin=4pt,
    aboveskip=8pt,
    belowskip=8pt
}

\newtheorem{conjecture}{Conjecture}

\title{Empirical Investigation of an Open Conjecture:\\
KLS Conjecture (Kannan–Lovász–Simonovits)}
\author{Agentic NL$\to$Lean~4 Pipeline \\ \small Job \#14}
\date{\today}

\begin{document}
\maketitle

\begin{abstract}
This report documents the empirical investigation of an open mathematical conjecture
that could not be formally proved or disproved in Lean~4 with Mathlib.
Numerical experiments were conducted to gather evidence for or against the conjecture.
The empirical verdict is: \textbf{\textcolor{green!70!black}{Empirically Supported}}.
The conjecture remains formally open.
\end{abstract}

\section{Conjecture Statement}

\begin{conjecture}
\begin{lstlisting}
KLS Conjecture (Kannan–Lovász–Simonovits)

For each integer 
𝑛
≥
1
n≥1, let 
𝜇
μ be a Borel probability measure on 
𝑅
𝑛
R
n
 that is absolutely continuous with respect to Lebesgue measure, with density 
𝑓
:
𝑅
𝑛
→
[
0
,
∞
)
f:R
n
→[0,∞) of the form 
𝑓
(
𝑥
)
=
𝑒
−
𝑉
(
𝑥
)
f(x)=e
−V(x)
, where 
𝑉
:
𝑅
𝑛
→
(
−
∞
,
∞
]
V:R
n
→(−∞,∞] is convex; equivalently, 
𝜇
μ is log-concave. Assume 
𝜇
μ is isotropic, meaning

∫
𝑅
𝑛
𝑥
 
𝑑
𝜇
(
𝑥
)
=
0
and
∫
𝑅
𝑛
𝑥
𝑥
⊤
 
𝑑
𝜇
(
𝑥
)
=
𝐼
𝑛
,
∫
R
n
	​

xdμ(x)=0and∫
R
n
	​

xx
⊤
dμ(x)=I
n
	​

,

where 
𝐼
𝑛
I
n
	​

 is the 
𝑛
×
𝑛
n×n identity matrix.

For a measurable set 
𝐴
⊆
𝑅
𝑛
A⊆R
n
, define its 
𝜀
ε-enlargement by 
𝐴
𝜀
=
{
𝑥
∈
𝑅
𝑛
:
dist
⁡
(
𝑥
,
𝐴
)
≤
𝜀
}
A
ε
	​

={x∈R
n
:dist(x,A)≤ε} and its Minkowski boundary measure (with respect to 
𝜇
μ) by

𝜇
+
(
𝐴
)
=
lim inf
⁡
𝜀
↓
0
𝜇
(
𝐴
𝜀
)
−
𝜇
(
𝐴
)
𝜀
.
μ
+
(A)=
ε↓0
liminf
	​

ε
μ(A
ε
	​

)−μ(A)
	​

.

Define the Cheeger (isoperimetric) constant of 
𝜇
μ as

𝜓
𝜇
=
inf
⁡
𝐴
 measurable
,
 
0
<
𝜇
(
𝐴
)
<
1
𝜇
+
(
𝐴
)
min
⁡
{
𝜇
(
𝐴
)
,
1
−
𝜇
(
𝐴
)
}
.
ψ
μ
	​

=
A measurable, 0<μ(A)<1
inf
	​

min{μ(A),1−μ(A)}
μ
+
(A)
	​

.
Unsolved Problem

(Kannan–Lovász–Simonovits conjecture; Kannan et al. (1995)) determine whether there exists a universal constant 
𝑐
>
0
c>0 (independent of 
𝑛
n and of 
𝜇
μ) such that for every dimension 
𝑛
n and every isotropic log-concave probability measure 
𝜇
μ on 
𝑅
𝑛
R
n
,

𝜓
𝜇
≥
𝑐
.
ψ
μ
	​

≥c.

Equivalently, if 
𝐶
𝑛
C
n
	​

 is the smallest number such that every isotropic log-concave 
𝜇
μ on 
𝑅
𝑛
R
n
 satisfies

𝜇
+
(
𝐴
)
≥
1
𝐶
𝑛
min
⁡
{
𝜇
(
𝐴
)
,
1
−
𝜇
(
𝐴
)
}
for all measurable 
𝐴
⊆
𝑅
𝑛
,
μ
+
(A)≥
C
n
	​

1
	​

min{μ(A),1−μ(A)}for all measurable A⊆R
n
,

the conjecture is that 
sup
⁡
𝑛
≥
1
𝐶
𝑛
<
∞
sup
n≥1
	​

C
n
	​

<∞ (that is, 
𝐶
𝑛
=
𝑂
(
1
)
C
n
	​

=O(1) as 
𝑛
→
∞
n→∞). Eldan (2013) introduced the stochastic localization approach. The best known general bound is currently 
𝐶
𝑛
=
𝑂
(
log
⁡
𝑛
)
C
n
	​

=O(
logn
	​

), due to Klartag (2023).

Solution Claims

Accepted clai
... [truncated]
\end{lstlisting}
\end{conjecture}

\section{Status}

\begin{center}
\fbox{\parbox{0.8\textwidth}{%
\textbf{Formal Status:} OPEN --- no Lean~4 proof or disproof was found.\\[4pt]
\textbf{Empirical Verdict:} \textcolor{green!70!black}{Empirically Supported}
}}
\end{center}

The pipeline attempted formal verification in Lean~4 with Mathlib but was unable to
produce a compiling proof or disproof. Empirical testing was then conducted to gather
numerical evidence.

\section{Basic Empirical Testing}

The following output was produced by the basic numerical experiment:

\begin{lstlisting}
=== EXPERIMENT PLAN ===

Goal: Empirically test the Kannan-Lovász-Simonovits (KLS) conjecture:
  inf_{n>=1} inf_{μ isotropic log-concave on R^n} ψ_μ  >  0,
where ψ_μ is the Cheeger (Minkowski-boundary-over-mass) constant.

Since we cannot enumerate ALL isotropic log-concave measures, we test a
diverse family of canonical examples across a broad range of dimensions:
  (a) Standard Gaussian                          -- ψ known = √(2/π)
  (b) Uniform on cube [-√3,√3]^n                 -- product, log-concave
  (c) Uniform on Euclidean ball  (radius √(n+2)) -- rotationally symmetric
  (d) Product Laplace (scale 1/√2)               -- heavy-tail log-concave
  (e) Uniform on ℓ1 ball                         -- non-product, convex body
Dimensions tested: n ∈ {2, 5, 10, 20, 50, 100, 200}; plus stress test n=500.

Methodology (multiple complementary angles):
  1) Half-space Cheeger estimate.  For each (μ, n) we draw N=3000 i.i.d.
     points, whiten them (subtract mean, multiply by Cov^{-1/2}) to enforce
     isotropy EXACTLY on the empirical measure, then for K=30 random unit
     directions u estimate
         R_u = inf_t  f_u(t) / min(F_u(t), 1 - F_u(t))
     via KDE + empirical CDF at quantiles q ∈ [0.03, 0.97].
     Set ψ̂_HS(n,μ) = min_u R_u (an upper bound on the true ψ_μ).
  2) Counterexample hunt by direction.  For the cube (believed extremal
     candidate) we evaluate axis and diagonal directions explicitly to
     see which directions try to "break" KLS.
  3) Non-half-space cuts.  Coordinate slabs {|x_1|≤t} and Euclidean
     balls {|x|≤r} are evaluated on Gaussian(n=50) to confirm that
     half-spaces are tighter (smaller ratio) than these alternatives.
  4) Asymptotic check.  We plot ψ̂_HS(n) and also ψ̂_HS(n)·√(log n).
     KLS predicts the former flat; Klartag (2023) upper bound is
     C_n=O(√log n), i.e. ψ̂·√log n ≥ const.  A rapid decay of ψ̂ toward
     0 would refute KLS.
  5) Statistical trend test.  Linear regression of ψ̂_HS on log(n) for
     each family; negative slopes of big magnitude would be a warning.

Verdict rule:
  - SUPPORTED if min over all (μ,n) tested of ψ̂_HS ≥ 0.30 AND the
    regression slopes vs log(n) are NOT sharply negative.
  - REFUTED if we find ψ̂ collapsing toward 0 in a dimension-growing way.
  - INCONCLUSIVE otherwise.

=== RUNNING: half-space Cheeger estimates ===
    n |         Gaussian |     Uniform cube |     Uniform ball |     Laplace prod |  Uniform ℓ1-ball
----------------------------------------------------------------------------------------------------
    2 |           0.7935 |           0.5736 |           0.6259 |           0.9676 |           0.5839
    5 |           0.7596 |           0.6052 |           0.6824 |           0.8484 |           0.6955
   10 |           0.7465 |           0.6961 |           0.7176 |           0.7931 |           0.6960
   20 |           0.7500 |           0.7470 |           0.7405 |           0.7850 |           0.7259
   50 |           0.7472 |           0.7564 |           0.7382 |           0.7857 |           0.7265
  100 |           0.7445 |           0.7553 |           0.7369 |           0.7477 |           0.7404
  200 |           0.7578 |           0.7343 |           0.7453 |           0.7685 |           0.7565

Total directional trials (half-space cuts): 1050
Theoretical Gaussian Cheeger: √(2/π) = 0.7979

=== BENCHMARK: non-half-space cuts (Gaussian, n=50) ===
  Coordinate slab {|x1|<=t}   min ratio = 1.2821
  Euclidean ball  {|x|<=r}   min ratio = 1.1598
  Half-space est at n=50                 = 0.7472
  (Half-space < slab/ball -> half-spaces are the tighter test sets, consistent with KLS expectations.)

=== COUNTEREXAMPLE HUNT: extremal directions in the cube (n=100) ===
  axis direction (worst?):  ψ̂ = 0.5807  (theory: uniform[-√3,√3] => 1/√3 ≈ 0.5774)
  diagonal direction (CLT): ψ̂ = 0.8032  (theory: ~Gaussian => √(2/π) ≈ 0.7979)
  random direction:         ψ̂ = 0.7675

=== HIGH-DIM STRESS TEST: n = 500 ===
          Gaussian, n=500: ψ
... [truncated]
\end{lstlisting}


\section{Experiment Code (Basic)}

\begin{lstlisting}[language=Python]
import numpy as np
import matplotlib
matplotlib.use("Agg")
import matplotlib.pyplot as plt
from scipy import stats
import time

np.random.seed(42)
rng = np.random.default_rng(42)

start = time.time()

print("=== EXPERIMENT PLAN ===")
print("""
Goal: Empirically test the Kannan-Lovász-Simonovits (KLS) conjecture:
  inf_{n>=1} inf_{μ isotropic log-concave on R^n} ψ_μ  >  0,
where ψ_μ is the Cheeger (Minkowski-boundary-over-mass) constant.

Since we cannot enumerate ALL isotropic log-concave measures, we test a
diverse family of canonical examples across a broad range of dimensions:
  (a) Standard Gaussian                          -- ψ known = √(2/π)
  (b) Uniform on cube [-√3,√3]^n                 -- product, log-concave
  (c) Uniform on Euclidean ball  (radius √(n+2)) -- rotationally symmetric
  (d) Product Laplace (scale 1/√2)               -- heavy-tail log-concave
  (e) Uniform on ℓ1 ball                         -- non-product, convex body
Dimensions tested: n ∈ {2, 5, 10, 20, 50, 100, 200}; plus stress test n=500.

Methodology (multiple complementary angles):
  1) Half-space Cheeger estimate.  For each (μ, n) we draw N=3000 i.i.d.
     points, whiten them (subtract mean, multiply by Cov^{-1/2}) to enforce
     isotropy EXACTLY on the empirical measure, then for K=30 random unit
     directions u estimate
         R_u = inf_t  f_u(t) / min(F_u(t), 1 - F_u(t))
     via KDE + empirical CDF at quantiles q ∈ [0.03, 0.97].
     Set ψ̂_HS(n,μ) = min_u R_u (an upper bound on the true ψ_μ).
  2) Counterexample hunt by direction.  For the cube (believed extremal
     candidate) we evaluate axis and diagonal directions explicitly to
     see which directions try to "break" KLS.
  3) Non-half-space cuts.  Coordinate slabs {|x_1|≤t} and Euclidean
     balls {|x|≤r} are evaluated on Gaussian(n=50) to confirm that
     half-spaces are tighter (smaller ratio) than these alternatives.
  4) Asymptotic check.  We plot ψ̂_HS(n) and also ψ̂_HS(n)·√(log n).
     KLS predicts the former flat; Klartag (2023) upper bound is
     C_n=O(√log n), i.e. ψ̂·√log n ≥ const.  A rapid decay of ψ̂ toward
     0 would refute KLS.
  5) Statistical trend test.  Linear regression of ψ̂_HS on log(n) for
     each family; negative slopes of big magnitude would be a warning.

Verdict rule:
  - SUPPORTED if min over all (μ,n) tested of ψ̂_HS ≥ 0.30 AND the
    regression slopes vs log(n) are NOT sharply negative.
  - REFUTED if we find ψ̂ collapsing toward 0 in a dimension-growing way.
  - INCONCLUSIVE otherwise.
""")

# ----------------------- Samplers ------------------------

def sample_gauss(N_, n):
    return rng.standard_normal((N_, n))

def sample_cube(N_, n):
    return rng.uniform(-np.sqrt(3.0), np.sqrt(3.0), (N_, n))

def sample_uniform_ball(N_, n):
    R = np.sqrt(n + 2.0)
    Y = rng.standard_normal((N_, n))
    Y /= np.linalg.norm(Y, axis=1, keepdims=True)
    u = rng.random(N_)
    r = R * u**(1.0/n)
    return r[:, None] * Y

def sample_laplace(N_, n):
    return rng.laplace(0.0, 1.0/np.sqrt(2.0), (N_, n))

def sample_uniform_l1_ball(N_, n):
    R = np.sqrt((n + 1.0) * (n + 2.0) / 2.0)
    E = rng.exponential(1.0, (N_, n + 1))
    S = E.sum(axis=1, keepdims=True)
    T = R * E[:, :n] / S
    signs = rng.choice(np.array([-1.0, 1.0]), size=(N_, n))
    return signs * T

distros = {
    "Gaussian":         sample_gauss,
    "Uniform cube":     sample_cube,
    "Uniform ball":     sample_uniform_ball,
    "Laplace prod":     sample_laplace,
    "Uniform ℓ1-ball":  sample_uniform_l1_ball,
}

# ----------------------- Helpers -------------------------

def isotropize(X):
    Xc = X - X.mean(axis=0)
    if X.shape[1] == 1:
        s = Xc.std(ddof=0)
        return Xc / max(s, 1e-12)
    C = np.cov(Xc, rowvar=False)
    vals, vecs = np.linalg.eigh(C)
    vals = np.maximum(vals, 1e-12)
    W = (vecs * vals**(-0.5)) @ vecs.T
    return Xc @ W

def halfspace_cheeger_estimate(X, n_dirs=30, q_lo=0.03, q_hi=0.97, n_q=60):
    N_, n = X.shape
    U = rng.standard_normal((n_dirs, n))
    U /= np.linalg.norm(U, axis=1, keepdims=True)
    Z = X @ U.T
    qs = np.linspace(q_lo, q_hi, n_q)
    ratios = np.zeros(n_dirs)
    for k in range(n_dirs):
        z = Z[:, k]
        kde = stats.gaussian_kde(z, bw_method="silverman")
        ts = np.quantile(z, qs)
        fs = kde(ts)
        rs = fs / np.minimum(qs, 1.0 - qs)
        ratios[k] = rs.min()
    return ratios.min(), ratios

# ----------------------- Main sweep ----------------------

dims = [2, 5, 10, 20, 50, 100, 200]
N_SAMPLES = 3000
N_DIRS = 30

results = {name: {"psi_min": [], "psi_mean": [], "psi_std": []}
           for name in distros}
total_trials = 0

print("=== RUNNING: half-space Cheeger estimates ===")
header = f"{'n':>5} | " + " | ".join(f"{name:>16}" for name in distros)
print(header)
print("-" * len(header))
for n in dims:
    row = [f"{n:>5}"]
    for name, sampler in distros.items():
        X = sampler(N_SAMPLES, n)
        X = isotropize(X)
        psi_hat, per_dir = halfspace_
# ... [truncated]
\end{lstlisting}


\section{Conclusion}

The conjecture remains formally open. Numerical experiments \textbf{support} the conjecture --- no counterexamples were found across all tested parameter ranges.
Further investigation (both formal and empirical) is warranted.

\end{document}
